\maketitle
%\thispagestyle{fancy}
\section*{Overview}

\begin{enumerate}
  \item \emph{Price stability.}
    We formulated quantitative price stability and predictability objectives and began seeking alignment with Swarm Research.
    %
    The strategy and timeline for improvements on price stability is still under development.

  \item \emph{Withdrawals and stake record updates.} 
    The SWIP is on hold pending a security review.

\end{enumerate}

\section*{Project management}

\subsection*{Completed}

\begin{itemize}
  \item Price stability: Research into history of cloud object storage pricing.
\end{itemize}

\subsection*{Ongoing}

\begin{itemize}
  \item Price stability: Build alignment around user-facing price stability objectives framework. 
  \item Price stability: Exploration of price controller design space.
  \item Price stability: Develop risk model for supply-side stability under pricing changes
\end{itemize}

\subsection*{Backlog}
\begin{itemize}
  \item Migration-free stake registry: triage and concrete proposal.
  \item Negative incentives: explore and evaluate solutions.
  \item Consensus attacks: analyse variants, evaluate threat level, list mitigations.
  \item Node dispersion: build mathematical model for address allocation schemes and node dispersion.
  \item Big nodes: Establish framework for converting findings into contributions back into mainline bee
\end{itemize}


\section*{Current work}

\subsection*{Stake withdrawals}

The stake withdrawals upgrade is currently on hold pending a security review.

\subsection*{Price stability}
%
We continue to structure the work in three workstreams:
\begin{enumerate}
  \item \emph{What do we want to happen?} Identify user and supplier needs, turn them into metrics, objectives, and thresholds.
  \item \emph{What has been happening?} Measure and analyse metrics from Swarm and the wider market.
  \item \emph{How do we get to our objectives?} Solution generation and evaluation.
\end{enumerate}
%
This month our focus was on setting up early candidates for objective setting and seeking preliminary alignment with Swarm Research.

\paragraph{What do we want to happen?}
%
It's essential to align on clear goals for price stability before executing on any solution strategy.
%
We presented some initial suggestions for price stability and predictability objectives to the Research Team, based on two principles:
%
\begin{enumerate}
  \item \emph{Short term billing certainty.} Swarm should be able to commit to a firm quote for storage over an upcoming accounting epoch, most likely the calendar month.
  \item \emph{Trust in long term price stability.} The month on month change in Swarm prices should be ``small'' except in exceptional circumstances.
\end{enumerate}
%
To make the second objective precise, we need to quantify ``small'' and ``exceptional circumstances''. A reasonable yardstick for ``small'' are FX swings between major fiat currencies — since Swarm is global in nature, the storage price will certainly be exposed to these at the very least.
%
There are many ways to specify ``exceptional circumstances,'' but as a starting point we propose measuring \emph{frequency} and \emph{attribution to known phenomena}.
%
Then we can specify and track objectives such as ``eliminate price swings due to client behaviour at network radius changes'' or ``reduce price swings that are not associated to a known risk indicator of Swarm supply to at most once per year.''

To provide evidence that these objectives are feasible, we'll do simulations and sketch approaches to implement variants of the objectives at different layers — see below.

\paragraph{What has been happening?}

See Figure \ref{fig:monthly-cost-estimate} for an estimate of the monthly dollar cost of Swarm storage over the past 12 months.
%
The difference between this cost sequence and the typical tradcloud experience is stark: 
\begin{itemize}
  \item The most striking phenomenon is the 150\% price jump in November, which we know was substantially influenced by a client behaviour that caused the reveal count to drop to zero as the network resident size dropped a notch.

  The price returned to its previous levels three months later but continued to drop. Arguably, the mechanism is still in price discovery — and we have reason to believe that node operators are currently running at a loss, suggesting that another large price increase may be needed in the near future.
  \item Even before November, where monthly prices were comparatively stable, July/August saw a price increase of over 20\%. At any tradcloud provider, this would be a newsworthy event; for Swarm, it's just another Tuesday. 
\end{itemize}

\begin{figure}[ht]
  \centering
  \includegraphics[width=\linewidth]{monthly-cost-estimate}
  \caption{Monthly cost estimate.}
  \label{fig:monthly-cost-estimate}
\end{figure}

\paragraph{How do we get to our objectives?}

Price stability solutions can live at multiple layers of Swarm's architecture: protocol, middleware, admin policy, client, application, or presentation, to name a few. 
%
There is currently interest in the Swarm research community in exploring non-protocol interventions:
%
\begin{itemize}
  \item \textbf{Policy.} Centralised price setting by the Foundation;
  \item \textbf{Middleware.} Third party intermediary offering fixed pricing and internalising Swarm price risk.
\end{itemize}
%
Either of these solutions could likely be deployed quickly, since they would not require going through the kind of rigorous evaluation that a protocol change would demand.

If we wish to explore in-protocol or policy solutions, we will need a risk model to measure and forecast the impact of changes on Swarm service reliability.
%
A risk model comprises two components:
%
\begin{itemize}
  \item \emph{What outcomes would we consider acceptable?} Acceptance thresholds for supply side outcomes, i.e.~node availability.
   
  We base our thresholds on two core objectives for Swarm service reliability:
  \begin{enumerate}
    \item \emph{Availability.} Minimise the rate $A$ of failed upload/download requests.
    \item \emph{Durability.} Minimise the rate $B$ of irrecoverably lost data objects.
  \end{enumerate}
  The threshold for $B$ will need to be much lower than the threshold for $A$.
  %
  Going by advertised reliability figures from tradcloud storage, a typical threshold for $A$ would be an API request failure (5xx error) rate of $10^{-3}$, while $B$ should sit at under $10^{-9}$ objects lost per year.

  Today, Swarm's availability and durability performance is upheld by the Foundation's not-for-profit node fleet.
  %
  However, ultimately Swarm will need to be able to function without not-for-profit operators; a risk model must therefore try to assess reliability based on the distribution of staked nodes.

  \item \emph{What might happen under the proposed changes?} To answer this question, we need a behavioural model of how operators will respond to protocol and market environment.

  For example, implementing central price setting by the Foundation requires the Foundation to choose a price that's enough to cover the costs of an adequately redundant node population.
  %
  The behavioural response we need to model here is the number of nodes the market will supply at a given price.
    
\end{itemize}